% GENERATED FILE. Edit canonical lessons, then run npm run generate:books.
% canonical-source-hash: fnv1a64:80cf86c58c261016
% canonical-lessons: SA-C283-apeksaya, SA-C283-arabhya-paryantam, SA-C283-idrsah, SA-S229-letter-ii, SA-R283-comparing-and-time

\chapter{Bigger, Better, and How Long}
\label{ch:sa-bigger-better-and-how-long}
\hlchaptermodality{283}

\begin{chapteropening}
\textbf{By the end of this chapter:} \emph{I can compare two things (nagaraṃ grāmāpekṣayā mahat; grāmāt mahattaram), say which is best (sarveṣu … śreṣṭhaḥ), and say how long something lasts (prātaḥ ārabhya sāyaṃ paryantam; ekāṃ horām).}

Complete SA-R283-comparing-and-time: compare two things, say which is best, and say how long something lasted.
\end{chapteropening}

\section[\texorpdfstring{\sk{अपेक्षया} (apekṣayā)}{अपेक्षया (apekṣayā)}]{\sk{अपेक्षया} (apekṣayā) — compared with}

\label{lesson:SA-C283-apeksaya}



\begin{quote}
Before the new one: say \textquotedblleft{}big\textquotedblright{}. (\emph{mahat}.) Then say \textquotedblleft{}the city is big\textquotedblright{}. (\emph{nagaraṃ mahat asti}.) Now set it beside the village.
\end{quote}

\subsection*{You'll want to know: Y-\sk{अपेक्षया}}
\begin{quote}
\textbf{\sk{नगरं} \sk{ग्रामापेक्षया} \sk{महत्।}}
\end{quote}

\begin{quote}
\emph{nagaraṃ grāmāpekṣayā mahat.} — \textquotedblleft{}The city is bigger than the village.\textquotedblright{}
\end{quote}

\textbf{\sk{अपेक्षया}} (\emph{apekṣayā}), \textquotedblleft{}compared with\textquotedblright{}, joins the end of the thing compared, as one word: \emph{grāma} and \emph{apekṣayā} meet as \textbf{\sk{ग्रामापेक्षया}} (\emph{grāmāpekṣayā}), two short \emph{a}s making one long \emph{ā}. The describing word stays as it is: \textbf{\sk{महत्}}, \textquotedblleft{}big\textquotedblright{}, says \textquotedblleft{}bigger\textquotedblright{} here.

\begin{quote}
\textbf{\sk{उद्यानं} \sk{वनापेक्षया} \sk{सुन्दरम्।}}
\end{quote}

\begin{quote}
\emph{udyānaṃ vanāpekṣayā sundaram.} — \textquotedblleft{}The garden is more beautiful than the forest.\textquotedblright{}
\end{quote}

\begin{quote}
\textbf{\sk{धूमयानं} \sk{लोकयानापेक्षया} \sk{शीघ्रं} \sk{गच्छति।}}
\end{quote}

\begin{quote}
\emph{dhūmayānaṃ lokayānāpekṣayā śīghraṃ gacchati.} — \textquotedblleft{}The train goes faster than the bus.\textquotedblright{}
\end{quote}

\subsection*{You'll want to know: -\sk{आत्} with -\sk{तर}, and the best}
The other way puts the thing compared in the \emph{-āt} form of \textbf{\sk{भारतात्}} (\emph{bhāratāt}), \textquotedblleft{}from India\textquotedblright{}, and adds \textbf{-\sk{तर}} (\emph{-tara}), \textquotedblleft{}more\textquotedblright{}, to the describing word (Pāṇini 2.3.42):

\begin{quote}
\textbf{\sk{नगरं} \sk{ग्रामात्} \sk{महत्तरम्।}}
\end{quote}

\begin{quote}
\emph{nagaraṃ grāmāt mahattaram.} — \textquotedblleft{}The city is bigger than the village.\textquotedblright{}
\end{quote}

\begin{quote}
\textbf{\sk{उद्यानं} \sk{वनात्} \sk{सुन्दरतरम्।}}
\end{quote}

\begin{quote}
\emph{udyānaṃ vanāt sundarataram.} — \textquotedblleft{}The garden is more beautiful than the forest.\textquotedblright{}
\end{quote}

For the best of all, \textbf{\sk{सर्वेषु}} (\emph{sarveṣu}), \textquotedblleft{}among all\textquotedblright{}, with \textbf{\sk{श्रेष्ठः}} (\emph{śreṣṭhaḥ}), \textquotedblleft{}best\textquotedblright{}:

\begin{quote}
\textbf{\sk{सर्वेषु} \sk{फलेषु} \sk{आम्रं} \sk{श्रेष्ठम्।}}
\end{quote}

\begin{quote}
\emph{sarveṣu phaleṣu āmraṃ śreṣṭham.} — \textquotedblleft{}Of all fruits, the mango is the best.\textquotedblright{}
\end{quote}

\begin{quote}
\textbf{\sk{सर्वेषु} \sk{मीरा} \sk{श्रेष्ठा।}}
\end{quote}

\begin{quote}
\emph{sarveṣu mīrā śreṣṭhā.} — \textquotedblleft{}Mira is the best of all.\textquotedblright{}
\end{quote}

\subsection*{Your turn}
Say these aloud:

\begin{itemize}
\raggedright
  \item \textquotedblleft{}The house is bigger than the shop\textquotedblright{} — \emph{gṛham āpaṇāpekṣayā mahat}
  \item \textquotedblleft{}The garden is bigger than the shop\textquotedblright{} — \emph{udyānam āpaṇāpekṣayā mahat}
  \item \textquotedblleft{}The city is bigger than the shop\textquotedblright{} — \emph{nagaram āpaṇāpekṣayā mahat}
\end{itemize}

\begin{tcolorbox}[breakable,colback=teal!4,colframe=teal!35!black,title={Before you move on}]
\textquotedblleft{}Compared with the village\textquotedblright{}? (\textbf{\sk{ग्रामापेक्षया}}.) \textquotedblleft{}Bigger than the village\textquotedblright{}, the other way? (\textbf{\sk{ग्रामात्} \sk{महत्तरम्}}.) \textquotedblleft{}The best\textquotedblright{}? (\textbf{\sk{श्रेष्ठः}}.)
\end{tcolorbox}
\section[\texorpdfstring{\sk{आरभ्य} · \sk{पर्यन्तम्} (ārabhya · paryantam)}{आरभ्य · पर्यन्तम् (ārabhya · paryantam)}]{\sk{आरभ्य} · \sk{पर्यन्तम्} (ārabhya · paryantam) — from, until, and how long}

\label{lesson:SA-C283-arabhya-paryantam}



\begin{quote}
Say \textquotedblleft{}morning\textquotedblright{} and \textquotedblleft{}evening\textquotedblright{}. (\emph{prātaḥ}, \emph{sāyam}.) Then say \textquotedblleft{}the city is bigger than the village\textquotedblright{}. (\emph{nagaraṃ grāmāpekṣayā mahat}.)
\end{quote}

\subsection*{You'll want to know: from … until}
\begin{quote}
\textbf{\sk{प्रातः} \sk{आरभ्य} \sk{सायं} \sk{पर्यन्तम्।}}
\end{quote}

\begin{quote}
\emph{prātaḥ ārabhya sāyaṃ paryantam.} — \textquotedblleft{}From morning until evening.\textquotedblright{}
\end{quote}

\textbf{\sk{आरभ्य}} (\emph{ārabhya}) is the \emph{-ya} form of \textbf{\sk{आरभते}} (\emph{ārabhate}), \textquotedblleft{}begins\textquotedblright{}: \textquotedblleft{}beginning from\textquotedblright{}, so \textquotedblleft{}from, since\textquotedblright{}. \textbf{\sk{पर्यन्तम्}} (\emph{paryantam}) is \textquotedblleft{}up to, until\textquotedblright{}.

\begin{quote}
\textbf{\sk{अहं} \sk{सायं} \sk{पर्यन्तं} \sk{कार्यं} \sk{करोमि।}}
\end{quote}

\begin{quote}
\emph{ahaṃ sāyaṃ paryantaṃ kāryaṃ karomi.} — \textquotedblleft{}I work until evening.\textquotedblright{}
\end{quote}

\begin{quote}
\textbf{\sk{सः} \sk{ह्यः} \sk{आरभ्य} \sk{रुग्णः} \sk{अस्ति।}}
\end{quote}

\begin{quote}
\emph{saḥ hyaḥ ārabhya rugṇaḥ asti.} — \textquotedblleft{}He has been ill since yesterday.\textquotedblright{}
\end{quote}

\subsection*{You'll want to know: for how long}
How long a thing lasts needs no word for \textquotedblleft{}for\textquotedblright{}: the time takes the \emph{-am} form, the form of the thing a verb acts on (Pāṇini 2.3.5). \textbf{\sk{होरा}} (\emph{horā}), \textquotedblleft{}an hour\textquotedblright{}, becomes \textbf{\sk{होराम्}}:

\begin{quote}
\textbf{\sk{अहम्} \sk{एकां} \sk{होरां} \sk{पठामि।}}
\end{quote}

\begin{quote}
\emph{aham ekāṃ horāṃ paṭhāmi.} — \textquotedblleft{}I read for an hour.\textquotedblright{}
\end{quote}

\begin{quote}
\textbf{\sk{सः} \sk{सर्वं} \sk{दिनं} \sk{कार्यं} \sk{करोति।}}
\end{quote}

\begin{quote}
\emph{saḥ sarvaṃ dinaṃ kāryaṃ karoti.} — \textquotedblleft{}He works the whole day.\textquotedblright{}
\end{quote}

For two, the dual: \textbf{\sk{द्वौ} \sk{मासौ}} (\emph{dvau māsau}), \textquotedblleft{}two months\textquotedblright{}. \textbf{\sk{स्थितवती}} (\emph{sthitavatī}) is the past of \textbf{\sk{तिष्ठति}}, \textquotedblleft{}stays\textquotedblright{}:

\begin{quote}
\textbf{\sk{सा} \sk{द्वौ} \sk{मासौ} \sk{ग्रामे} \sk{स्थितवती।}}
\end{quote}

\begin{quote}
\emph{sā dvau māsau grāme sthitavatī.} — \textquotedblleft{}She stayed in the village for two months.\textquotedblright{}
\end{quote}

\subsection*{Your turn}
Say these aloud:

\begin{itemize}
\raggedright
  \item \textquotedblleft{}I read for an hour\textquotedblright{} — \emph{aham ekāṃ horāṃ paṭhāmi}
  \item \textquotedblleft{}I read the whole day\textquotedblright{} — \emph{ahaṃ sarvaṃ dinaṃ paṭhāmi}
  \item \textquotedblleft{}I read until evening\textquotedblright{} — \emph{ahaṃ sāyaṃ paryantaṃ paṭhāmi}
\end{itemize}

\begin{tcolorbox}[breakable,colback=teal!4,colframe=teal!35!black,title={Before you move on}]
\textquotedblleft{}From morning\textquotedblright{}? (\textbf{\sk{प्रातः} \sk{आरभ्य}}.) \textquotedblleft{}Until evening\textquotedblright{}? (\textbf{\sk{सायं} \sk{पर्यन्तम्}}.) \textquotedblleft{}For an hour\textquotedblright{}? (\textbf{\sk{एकां} \sk{होराम्}}.) Which word means \textquotedblleft{}for\textquotedblright{}? (None: the time says it by its form.)
\end{tcolorbox}
\section[\texorpdfstring{\sk{ईदृशः} (īdṛśaḥ)}{ईदृशः (īdṛśaḥ)}]{\sk{ईदृशः} (īdṛśaḥ) — such, of this kind}

\label{lesson:SA-C283-idrsah}



\begin{quote}
Before the new one: say \textquotedblleft{}the mango is the best of all fruits\textquotedblright{}. (\emph{sarveṣu phaleṣu āmraṃ śreṣṭham}.)
\end{quote}

\subsection*{You'll want to know: \sk{ईदृशः}}
\textbf{\sk{ईदृशः}} — \emph{īdṛśaḥ} — \textquotedblleft{}such, of this kind\textquotedblright{}. It agrees with its noun: \textbf{\sk{ईदृशम्}} (\emph{īdṛśam}) for a neuter noun, \textbf{\sk{ईदृशी}} (\emph{īdṛśī}) for a feminine one.

\begin{quote}
\textbf{\sk{ईदृशः} \sk{अनुभवः} \sk{मह्यं} \sk{रोचते।}}
\end{quote}

\begin{quote}
\emph{īdṛśaḥ anubhavaḥ mahyaṃ rocate.} — \textquotedblleft{}I like an experience like this.\textquotedblright{}
\end{quote}

\begin{quote}
\textbf{\sk{ईदृशं} \sk{नगरम्} \sk{अहं} \sk{न} \sk{दृष्टवान्।}}
\end{quote}

\begin{quote}
\emph{īdṛśaṃ nagaram ahaṃ na dṛṣṭavān.} — \textquotedblleft{}I have not seen a city like this.\textquotedblright{}
\end{quote}

\begin{cousinweb}
\textbf{\sk{ईदृशः}} (\emph{īdṛśaḥ}) is built from \emph{ī}, \textquotedblleft{}this\textquotedblright{}, and the root √dṛś, \textquotedblleft{}to see\textquotedblright{}, the root of \textbf{\sk{द्रक्ष्यति}}: \textquotedblleft{}looking like this\textquotedblright{} (CDIAL 1611 gives the two parts). Monier-Williams compares Latin \emph{tālis}, \textquotedblleft{}such\textquotedblright{}. It lives on as Hindi \emph{aisā}, \textquotedblleft{}such, like this\textquotedblright{} (CDIAL 1611).
\end{cousinweb}

\subsection*{Your turn}
Say these aloud:

\begin{itemize}
\raggedright
  \item \emph{īdṛśaḥ}
  \item \textquotedblleft{}Such an experience\textquotedblright{} — \emph{īdṛśaḥ anubhavaḥ}
  \item \textquotedblleft{}Such a city\textquotedblright{} — \emph{īdṛśaṃ nagaram}
\end{itemize}

\begin{tcolorbox}[breakable,colback=teal!4,colframe=teal!35!black,title={Before you move on}]
What does \textbf{\sk{ईदृशः}} mean? (\textquotedblleft{}Such, of this kind\textquotedblright{}.) \textquotedblleft{}Such a city\textquotedblright{}? (\textbf{\sk{ईदृशं} \sk{नगरम्}}.)
\end{tcolorbox}
\section[\texorpdfstring{\sk{ई} (ī)}{ई (ī)}]{\sk{ई} — the long ī, standing alone}

\label{lesson:SA-S229-letter-ii}



\begin{quote}
Before the new one: say \textbf{\sk{ईदृशः}} (\emph{īdṛśaḥ}), \textquotedblleft{}such\textquotedblright{}. Then name the sign \textbf{\sk{ी}} and the letter \textbf{\sk{इ}}. (\emph{ī}, \emph{i}.)

The long sign \textbf{\sk{ी}} has a standing form, and \textbf{\sk{ईदृशः}} opens with it.
\end{quote}

\subsection*{Script you'll notice: \sk{ई}}
\textbf{\sk{ई}} — \emph{ī}.

It is an \textbf{independent vowel}: \textbf{\sk{इ}} (\emph{i}) made long. On a consonant the same vowel is the sign \textbf{\sk{ी}}, as in \textbf{\sk{नदी}} (\emph{nadī}).

What it is made of:

\begin{itemize}
\raggedright
  \item the continuous body of \sk{इ}
  \item a separate upper curl
  \item the top shirorekhā
\end{itemize}

You have met it here:

\begin{itemize}
\raggedright
  \item \textbf{\sk{ईदृशः}} \emph{īdṛśaḥ} — such
\end{itemize}

\subsection*{Writing: \sk{ई}}
\hlblockfigure{\detokenize{figures/SA-S229-letter-ii-filmstrip.pdf}}{How \sk{ई} is written, stroke by stroke}

\begin{itemize}
\raggedright
  \item \textbf{1.} descend the upright, turn through both bowls, and finish down-right through the tail without lifting
  \item \textbf{2.} lift at the headline junction and sweep the upper curl upward, left, and around to the right
  \item \textbf{3.} lift and draw the top shirorekhā left-to-right
\end{itemize}

\textbf{Pen lifts: 2.} The pen comes up 2 times and no more.

\begin{quote}
Verified modern printed teaching form; everyday handwriting may simplify the curves.
\end{quote}

\begin{quote}
This is one attested teaching order and not a national standard — handwriting here is taught with school-to-school variation. Source: Saurmandal, ‘Devanagari \sk{ई} stroke order.svg’, panels 1–3, Wikimedia Commons, 5 August 2023.
\end{quote}

\subsection*{Your turn}
\begin{itemize}
\raggedright
  \item \emph{Look:} at this, and find \sk{ई} in it
\end{itemize}

\begin{quote}
\sk{ईदृशः}
\end{quote}

\begin{itemize}
\raggedright
  \item \emph{Trace it:} \sk{ई} three times, saying \emph{ī} as you finish each one
  \item \emph{Look:} at \textbf{\sk{ईदृशं} \sk{नगरम्}} and find \sk{ई} once more
\end{itemize}

\begin{tcolorbox}[breakable,colback=teal!4,colframe=teal!35!black,title={Before you move on}]
Which character is this — \sk{ई}? What sound does it carry? (\emph{ī}.) How is it unlike \textbf{\sk{इ}}? (One more curl above, and a long sound.) Which word did it come from? (\textbf{\sk{ईदृशः}}, \textquotedblleft{}such\textquotedblright{}.)
\end{tcolorbox}
\section[(review)]{Review — bigger, better, and how long}

\label{lesson:SA-R283-comparing-and-time}



\begin{quote}
Say \textquotedblleft{}I am pleased\textquotedblright{}, then \textquotedblleft{}my dream is that I will become a doctor\textquotedblright{}. (\emph{ahaṃ prasannaḥ asmi}; \emph{mama svapnaḥ asti, yat ahaṃ vaidyaḥ bhaviṣyāmi}.)
\end{quote}

\subsection*{Your turn}
\begin{itemize}
\raggedright
  \item \emph{Say it:} The city is bigger than the village.
  \item \emph{Say it:} The train goes faster than the bus.
  \item \emph{Say it:} Of all fruits, the mango is the best.
  \item \emph{Say it:} I worked from morning until evening.
  \item \emph{Say it:} She stayed in the village for two months. I like such a life.
  \item \emph{Read it:} \textbf{\sk{ईदृशः}} aloud, and say the sign that is its sitting form
\end{itemize}

\begin{tcolorbox}[breakable,colback=teal!4,colframe=teal!35!black,title={Before you move on}]
\textbf{\sk{ग्रामापेक्षया}} or \textbf{\sk{ग्रामात्}}, with \textbf{\sk{महत्तरम्}}? (\textbf{\sk{ग्रामात्}}.) \textbf{\sk{एकां} \sk{होराम्}}: one hour, or the hour itself? (For one hour.) Now the other way round: \textquotedblleft{}until\textquotedblright{}? (\textbf{\sk{पर्यन्तम्}}.)
\end{tcolorbox}
